
\pagestyle{empty}
\section*{Lieber Frank,}
Dieses Dokument ist eine Zusammenfassung und Kollektion meines bisherigen Vorankommens. Ich beginne mit einer Einführung in die Diffenentialgeometrie, welche für dich nichts neues enthalten wird. Danach kommt eine kurze Einführung in die Morse Theorie. Hier ist denke ich folgendes für dich spannend: 
\begin{enumerate}
	\item Theorem \ref{thm: Simultaneous Diagonalisation of Quadratic Forms}. Hier führe ich den \textbf{refined Morse chart} ein, welcher neben der quadratischen Form von $f$ auch die Metrik im Ursprung diagonalisiert. 
	\item Definition \ref{def: Morse homology} definiert den Randoperator und danach definiere ich noch die Kohomologie.
\end{enumerate}
Kapitel 4 ist denke ich im Gesamten spannend. Dort führe ich zunächst die verwendete Filtrierung ein und dann die Spektralsequenz. Der Beweis von Theorem \ref{thm commuting connecting homomorphism} ist wohl der Höhepunkt. Ich habe auch zur Wiederholung in dem "Beweisidee" Container die Definition des verbindenden Homomorphismus der K-Theorie nochmal wiederholt. (Eine Einführung in die K-Theorie habe ich der Länge wegen weggelassen.) Der Beweis ist ziemlich technisch, hat zu diesem Zeitpunkt noch viele Lemma und ist auch noch nicht ganz fertig. Es fehlt noch die technisch saubere Ausarbeitung des letzten Arguments(das habe ich auch markiert). Hier wirst du denke ich viel wiederfinden von unseren Gesprächen.

Generell ist das hier mein \textit{Arbeitsdokument} in etwas angepasster Form. Daher ist leider \textbf{noch} nicht alles einheitlich und etwas durcheinander und die hübschen Bilder fehlen noch. Auch möchte ich dich warnen, dass ich es noch nicht auf Rechtschreibung korrigiert habe und dank meiner Rechtschreibschwäche daher wohl noch viele Fehler zu finden sind.
\newpage
\pagestyle{scrheadings}